\section{Introduction}
\label{sec:intro}

The \emph{Capacitated Facility Location Problem} (CFLP) is a classical model for
selecting a set of facilities and allocating customer demand to them under
capacity restrictions.
It appears in distribution network design, supply chain planning, emergency
service location, and telecommunications infrastructure
design~\citep{klose2005facility,daskin2011network,laporte2019location}.
In the \emph{Multi-Source} CFLP (MSCFLP), a customer's demand may be split
among several open facilities.
In the \emph{Single-Source} CFLP (SSCFLP), each customer must be assigned to one
facility; this additional assignment restriction substantially changes the
structure and computational difficulty of the problem~\citep{klose2005facility}.
Both variants are NP-hard~\citep{sridharan1995capacitated}.

This paper focuses on the large-scale MSCFLP.
Avella et al.~\citep{avella2008effective} established a widely used
computational benchmark for this setting and proposed a Lagrangean-relaxation
and core-problem heuristic for instances with millions of variables.
Guastaroba and Speranza~\citep{guastaroba2012kernel} introduced a
Kernel Search (KS) matheuristic for the same problem class
(which they refer to simply as the CFLP; we retain the MSCFLP label to
distinguish it from the single-source variant).
Their method solves the LP relaxation, identifies an initial kernel of
LP-open and LP-fractional facilities, and solves a sequence of restricted MILPs
that add promising LP-closed facilities through capacity-expanding buckets.
The approach is effective because the strong MSCFLP formulation includes
variable upper-bound constraints, which make the LP solution a reliable predictor
of the facilities and arcs used in high-quality integer solutions.
Speed controls in the original paper act on the restricted-MILP phase: bucket
count limits, variable-count caps per subproblem, a kernel-pruning rule based on
repeated non-selection, and a forcing constraint requiring each bucket subproblem
to use a facility from the current bucket.
A faster variant, B-KS(0), searches only over the initial LP-selected kernel.
The original paper frames KS performance as a balance between restricted-search
quality and the time required to solve the resulting MILPs.

On the exact-method side, Görtz and Klose~\citep{gortz2012lagrangian} proposed a
Lagrangian branch-and-bound that is competitive with the best exact approaches on
the smaller GK instances.
Fischetti, Ljubić, and Sinnl~\citep{fischetti2016benders} later showed that a
modern Benders decomposition embedded in a branch-and-cut solver proves optimality
for 210 out of 445 large-scale p* instances within a 50{,}000-second time limit,
and consistently produces smaller gaps than KS across all three test beds.
Their computational study confirms that solving the LP at the root node
dominates the total runtime for the largest instances: more than 90\% of overall
computing time may be spent in LP-based cut separation.
To the best of our knowledge, these are the only results available for exact
methods on the large-scale multi-source CFLP benchmark of Avella et al.

This paper takes a different starting point: rather than proposing a new
acceleration strategy, we ask whether the design choices embedded in G\&S~2012
are each well-founded.
A close reading of the method reveals three design choices, each of which can
be refined independently.

\textbf{LP construction cost.}
The full LP relaxation contains $|\J||\I|$ shipping variables.
At optimality, most carry positive reduced cost and contribute nothing to the
LP solution or to the KS guidance it produces.
Constructing the full matrix is therefore unnecessary: column generation recovers
the same bound, facility partition, and reduced-cost signal without ever
materialising the irrelevant variables.

\textbf{Arc selection bias.}
G\&S~2012 includes arc $(j,i)$ whenever its reduced cost falls below a global
threshold $\gamma$, the mean reduced cost over LP-active arcs.
The threshold provides no guarantee on subproblem size and is structurally
biased: a bucket facility $j \in \Jzero$ has LP dual $\alpha^*_j = 0$, so its
arc reduced cost $c_{ji} - \pi^*_i$ contains no capacity signal.
A kernel facility with positive dual receives a capacity-dual discount that
lowers its reduced costs; the same $\gamma$ therefore admits more kernel arcs
than bucket arcs, giving bucket facilities a systematically smaller arc set
in every restricted MILP.

\textbf{Cold-start stage transitions.}
G\&S~2012 passes no solution information between consecutive restricted MILPs.
When the arc set is modified between stages, the incumbent from the previous
subproblem may no longer be feasible: arcs it uses may simply be absent from the
new edge set.
Because feasibility cannot be guaranteed, the incumbent cannot be supplied as
a warm start, and the solver must rediscover a feasible solution from scratch at
every stage.
This is not an obvious limitation of warm starting in general; it is a
structural consequence of arc-set transitions that discard incumbent arcs.

We address each in turn.
The LP is solved by column generation.
Arc selection is replaced by a client-centric flat-$k$ criterion: for each
client, the $k$ kernel facilities with lowest reduced cost are included,
giving an exact subproblem-size guarantee and treating all facilities uniformly
regardless of their LP dual value.
Stage transitions are made feasibility-preserving by including every arc used by
the current incumbent when building the next stage's edge set; we prove that
this construction guarantees the incumbent is feasible in the next restricted
MILP, making it a valid warm start by construction.
The resulting repaired variant is \emph{Funnel Kernel Search} (FKS).
It retains the full architecture of G\&S~2012 --- the strong MSCFLP formulation,
the LP-defined kernel, reduced-cost guidance, and bucket expansion --- and
modifies only the three components identified above.

The contributions are as follows.
\begin{itemize}
  \item \textbf{Design analysis.}
        We examine the G\&S~2012 design and identify three aspects where
        targeted modifications yield consistent improvements:
        (1)~the LP is built in full when column generation recovers the same
        information at lower cost;
        (2)~the global $\gamma$ threshold for arc selection admits no size
        guarantee and treats bucket facilities less favorably, since their LP
        dual is zero and their arc reduced costs carry no capacity signal,
        causing the same $\gamma$ to admit fewer arcs than for kernel
        facilities;
        (3)~the restricted-MILP phase starts cold at each stage, forcing the
        solver to rediscover feasibility after every arc-set transition.
  \item \textbf{Fix~1 --- LP cost.}
        We apply column generation to the MSCFLP LP relaxation,
        retaining the bound, facility partition, and reduced-cost signal
        required by KS while avoiding full-matrix construction.
  \item \textbf{Fix~3 --- feasibility-preserving stage transitions.}
        We prove that including every incumbent arc in the next stage's edge set
        guarantees the previous solution is feasible in the next restricted MILP
        (Proposition~\ref{prop:feasibility}).
        This is a structural result, not a heuristic: feasibility is guaranteed
        by construction and makes the incumbent a provably valid MIP start at
        every stage.
  \item \textbf{Fix~2 --- arc selection.}
        With the feasibility guarantee in place, we introduce a client-centric
        flat-$k$ criterion that replaces the global $\gamma$ threshold.
        The flat-$k$ rule narrows the fill set aggressively across stages
        (50$\to$20$\to$10 arcs per client) without risking cold restarts,
        because Proposition~\ref{prop:feasibility} holds regardless of how
        small $k$ becomes.
        It also gives an exact subproblem-size guarantee
        (Proposition~\ref{prop:size}) and treats all facilities uniformly.
  \item \textbf{Component-wise study.}
        The computational study isolates each fix via a dedicated comparison:
        full LP versus CG LP (Fix~1 alone), KS-CG versus KS-CG-WS (Fix~3
        alone), KS-CG-WS versus FKS-CG-WS (Fix~2 alone), and end-to-end
        KS-full versus FKS-CG-WS (combined effect).
\end{itemize}

The result is a matheuristic that is not only faster than the original but
also consistently finds solutions of higher quality. On the 110 Test Bed~A
instances with known proven optima, FKS finds the optimal solution on 109.

The remainder of the paper is organized as follows.
Section~\ref{sec:formulation} gives the MSCFLP formulation and the LP structure
used by KS.
Section~\ref{sec:algorithm} examines these three design aspects formally,
then presents the modifications in order of dependency: column-generation LP
(Fix~1), feasibility-preserving stage transitions with their structural proof
(Fix~3), and flat-$k$ arc selection (Fix~2) --- which relies on Fix~3 to
narrow arc sets safely.
Section~\ref{sec:experiments} describes the component-wise computational study
and reports results on the large-scale Test Bed benchmarks.
Section~\ref{sec:conclusion} concludes.
